\answer{ex:08:1}{(a)~$\approx1.7\cdot10^6$. (b)~$g=\frac{0.9}{\sqrt{0.19}}\,\sigma_{\text{después}}/\sigma_{\text{antes}}\approx16.5$: la respuesta solo depende del cociente de los ruidos.}
\answer{ex:08:2}{(a)~Valores propios $-(1\pm g\beta)/\tau$; la diferencia se amortigua en $\tau/(1-g\beta)$: $40\ms$ con $g\beta=0.5$ (la diferencia de entradas amplificada tres veces), $200\ms$ con $0.9$ ($19$ veces). (b)~$0.905$ y $1.105$; entre ambos límites solo gana la entrada fuerte.}
\answer{ex:08:3}{$P(k)=e^{gu_k/\sigma}/\sum_le^{gu_l/\sigma}$, es decir,~\eqref{eq:softmax}; $g=10\ln9\approx22$.}
\answer{ex:08:4}{$\bar\Delta=1/3$, $g^*\approx3\ln(1/h)$: $21$, $14$, $9$ frente a $16$--$23$, $11$ y $8$ en el modelo.}
\answer{ex:08:5}{$g^*$ de $16$--$23$ con $h=0.001$ a $6$--$8$ con $h=0.2$; la estimación $3\ln(1/h)$ vale dentro de un factor de uno y medio con $h\le0.05$; con $h\gtrsim\eta/10$ la regla $Q\leftarrow Q+\eta(R-Q)$ no llega a asimilar lo obtenido explorando, y $g^*$ se estanca en torno a $8$.}
\answer{ex:08:6}{$T_p=\tau\ln9=44\ms$ con $g_{\text{lo}}=0$ (modelo, $44\ms$) y $70\ms$ con $g_{\text{lo}}=0.25$: una ráfaga de $2$--$3\tau$ con ganancia cercana a cero.}
\answer{ex:08:7}{(a)~Mejor constante, $0.925$ ($g=8$); oráculo, $0.929$ ($16/8$): casi no hay nada que ganar. (b)~Por ejemplo, épocas tranquilas con recompensas escasas ($p\in[0,0.3]$) y cambiantes con $h=0.1$: $0.850$ frente a $0.872$; el ajuste obtiene cerca de una cuarta parte con $\eta=0.2$ y casi todo con $\eta=0.05$.}
